\section{Fixed moments do not give a polynomial rate}
\label{p1:moment-obstruction}

This example isolates the logical role of the growing-order
estimate. It is a family of probability laws, with no claimed
multiplicative realization.

\begin{proposition}[A perturbation invisible at each fixed order]
\label{p1:obstruction:fixed}
For \(0\le a\le1\) put
\[
 Q_a=\Pi_1+\frac a2\Pi_1*\Delta^{*2}.
\]
For every integer \(m\ge3\), define
\[
 \epsilon_{m,a}=\frac{(1-a/2)e^{-1}}{2m!},
 \qquad
 P_{m,a}=Q_a+\epsilon_{m,a}\Delta^{*m}.
\]
Then \(P_{m,a}\) is a probability measure satisfying
\(\frac12Q_a\le P_{m,a}\le\frac32Q_a\) pointwise.
Its factorial moments agree exactly with those of \(Q_a\)
at every order other than \(m\), and
\[
 \TV(P_{m,a},Q_a)
       =\frac{(1-a/2)2^{m-2}}{e\,m!}.
\]
For every fixed \(t\ge0\), its exponential moment
\(\E_{P_{m,a}}e^{tX}\) is bounded uniformly in \(m,a\).
\end{proposition}

\begin{proof}
The second Charlier polynomial at mean one is
\(k^2-3k+1\), whose minimum on \(\N\) is \(-1\).
Thus
\[
 Q_a\{k\}\ge(1-a/2)e^{-1}/k!.
\]
For \(0\le k\le m\),
\[
 \frac{\epsilon_{m,a}\binom mk}{Q_a\{k\}}
       \le\frac1{2(m-k)!}\le\frac12 .
\]
Outside this range the perturbation is zero. Positivity
and the two pointwise bounds follow. The perturbation
has total mass zero, so \(P_{m,a}\) has mass one.

The probability generating function of \(\Delta^{*m}\)
is \((z-1)^m\). Its derivatives at \(z=1\) vanish at
every order other than \(m\), where the derivative is
\(m!\). This proves the factorial-moment statement.
Its \(\ell^1\) norm is \(2^m\), yielding the exact
total variation distance. Finally,
\[
 \E_{P_{m,a}}e^{tX}
 \le\frac32
       \exp(e^t-1)\left(1+\frac12(e^t-1)^2\right),
\]
which is independent of \(m,a\).
\end{proof}

Take \(m_N=\lfloor\sqrt{\log N}\rfloor\) and
\(a_N^\circ=(\log N)N^{-2/3}\), for \(N\) sufficiently
large. At every fixed order \(r\), the factorial
moments of \(P_{m_N,a_N^\circ}\) eventually equal
\(1+\binom r2a_N^\circ\) exactly. Nevertheless,
Stirling's formula gives
\[
 \log\TV(P_{m_N,a_N^\circ},Q_{a_N^\circ})
   =-\tfrac12\sqrt{\log N}\log\log N
                   +O(\sqrt{\log N}).
\]
This distance exceeds \(N^{-\gamma}\) for every fixed
\(\gamma>0\), eventually. Fixed-order information,
even with the displayed exponential-moment control,
therefore cannot replace the uniform estimate on the
boundary of the stopped expansion.
